% =====================================================================
\section{Светодиод: закон Ома и мигание}
\label{les:led}
% =====================================================================
\lessonintro
  {правильно подключить светодиод, рассчитать резистор по закону Ома, заставить светодиод мигать}
  {урок~\ref{les:signals} (цифровой сигнал), урок~\ref{les:pins} (выводы GPIO), урок~\ref{les:ide} (загрузка программы)}
  {мигающий светодиод и «бегущий огонь» из четырёх светодиодов}
  {макетная плата, провода, 4 светодиода, 4 резистора \SI{220}{\ohm}}

\subsection{Светодиод}

\textbf{Светодиод} (англ. LED --- \emph{light-emitting diode}) --- это деталь, которая светится, когда через
неё течёт ток. Он пропускает ток \emph{только в одну сторону}, как дверь, которая открывается
только наружу. Поэтому у светодиода есть «плюс» и «минус»:

\begin{figure}[H]
\centering
\begin{tikzpicture}[x=1cm,y=1cm]
  % корпус
  \fill[red!75,opacity=0.85] (-0.6,1.2) -- (-0.6,2.6) arc[start angle=180,end angle=0,radius=0.6] -- (0.6,1.2) -- cycle;
  \fill[red!55!black] (-0.75,1.0) rectangle (0.62,1.2);
  \fill[white,opacity=0.5] (-0.3,2.7) circle (0.13);
  \draw[very thick,black!55] (-0.25,1.0) -- (-0.25,-1.2);
  \draw[very thick,black!55] (0.25,1.0) -- (0.25,-0.6);
  \node[font=\sffamily\small,align=right,anchor=east] at (-0.8,-0.9) {длинная ножка ---\\ \textbf{анод} («+»)};
  \node[font=\sffamily\small,align=left,anchor=west] at (0.8,-0.4) {короткая ножка ---\\ \textbf{катод} («−»)};
  \draw[thin,gray] (0.62,1.1) -- (1.3,1.5) node[right,font=\sffamily\small,text=black,align=left]{срез на ободке\\ со стороны катода};
  % обозначение
  \begin{scope}[xshift=7cm,yshift=0.8cm]
    \draw (0,0) node[left,font=\sffamily\small]{анод (+)} to[led,o-o] (3,0) node[right,font=\sffamily\small]{катод (−)};
    \draw[->,thick,esp] (0.9,-0.7) -- node[below,font=\sffamily\scriptsize]{ток течёт так} (2.1,-0.7);
    \node[font=\sffamily\small] at (1.5,1.2) {обозначение на схемах:};
  \end{scope}
\end{tikzpicture}
\caption{Светодиод и его условное обозначение. Треугольник показывает, в какую сторону течёт ток}
\end{figure}

Если перепутать ножки, светодиод просто не загорится --- ничего страшного. А вот если подключить
его без резистора, он сгорит. Разберёмся почему.

\subsection{Зачем нужен резистор: закон Ома}

Электричество удобно сравнивать с водой в трубах:
\begin{itemize}
  \item \textbf{напряжение} $U$ (вольты, В) --- это «давление», которое толкает воду;
  \item \textbf{ток} $I$ (амперы, А) --- сколько воды протекает за секунду;
  \item \textbf{сопротивление} $R$ (омы, \si{\ohm}) --- насколько узкая труба: чем уже, тем меньше воды пройдёт.
\end{itemize}

\begin{figure}[H]
\centering
\begin{tikzpicture}[x=1cm,y=1cm]
  \fill[blue!12] (0,0) rectangle (1.6,3);
  \fill[blue!35] (0,0) rectangle (1.6,2.3);
  \draw[thick] (0,3.2) -- (0,0) -- (1.6,0) -- (1.6,3.2);
  \node[font=\sffamily\small,align=center] at (0.8,3.6) {давление\\ = напряжение $U$};
  \fill[blue!35] (1.6,0.1) rectangle (3.6,0.6);
  \fill[blue!35] (3.6,0.25) rectangle (5.4,0.45);
  \fill[blue!35] (5.4,0.1) rectangle (7.6,0.6);
  \draw[thick] (1.6,0.6) -- (3.6,0.6) -- (3.6,0.45) -- (5.4,0.45) -- (5.4,0.6) -- (7.6,0.6);
  \draw[thick] (1.6,0.1) -- (3.6,0.1) -- (3.6,0.25) -- (5.4,0.25) -- (5.4,0.1) -- (7.6,0.1);
  \node[font=\sffamily\small,align=center] at (4.5,1.4) {узкое место\\ = сопротивление $R$};
  \draw[->,very thick,esp] (6,-0.3) -- (7.4,-0.3) node[right,font=\sffamily\small,text=black]{поток = ток $I$};
\end{tikzpicture}
\caption{«Водяная» модель электрической цепи}
\end{figure}

Эти три величины связаны \textbf{законом Ома} --- главным законом электроники:
\[
  I = \frac{U}{R}
  \qquad\text{или}\qquad
  R = \frac{U}{I}.
\]
Больше давление --- больше поток. Уже труба --- меньше поток.

\subsubsection{Почему светодиод сгорает без резистора}

Светодиод ведёт себя необычно: пока напряжение на нём маленькое, ток почти не течёт. Но стоит
напряжению превысить определённый порог --- \textbf{прямое напряжение} $U_F$, для красного светодиода
около 2~В, --- ток начинает расти \emph{очень} быстро:

\begin{figure}[H]
\centering
\begin{tikzpicture}
\begin{axis}[
  width=0.62\textwidth, height=5.4cm,
  xlabel={напряжение на светодиоде, В}, ylabel={ток, мА},
  xmin=0, xmax=2.4, ymin=0, ymax=60,
  grid=major, grid style={gray!25},
  tick label style={font=\scriptsize}, label style={font=\small},
]
  \addplot[very thick,esp,samples=150,domain=0:2.2] {min(80,0.5*exp((x-1.75)*14))};
  \addplot[dashed,okgreen!70!black,thick] coordinates {(0,6) (2.4,6)};
  \node[font=\scriptsize,anchor=south west,text=okgreen!50!black] at (axis cs:0.05,6) {нужный ток $\approx$ 6 мА};
  \node[circle,fill=okgreen!70!black,inner sep=1.5pt] at (axis cs:1.93,6) {};
  \draw[<-,thin] (axis cs:2.08,50) -- (axis cs:1.2,45) node[left,font=\scriptsize,align=right]{чуть больше напряжения ---\\ ток огромный,\\ светодиод сгорает};
  \draw[<-,thin] (axis cs:1.0,0.8) -- (axis cs:0.7,20) node[above,font=\scriptsize,align=center]{ток почти\\ не течёт};
\end{axis}
\end{tikzpicture}
\caption{Как ток через красный светодиод зависит от напряжения на нём (вольт-амперная характеристика)}
\end{figure}

Вывод платы выдаёт 3,3~В --- намного больше, чем нужно светодиоду. Если подключить его напрямую, ток
окажется огромным, и светодиод (а может, и вывод платы) сгорит. \textbf{Резистор} забирает на себя
«лишнее» напряжение и ограничивает ток. Поэтому правило:

\begin{center}
\fcolorbox{esp}{esp!6}{\parbox{0.8\textwidth}{\centering\sffamily\bfseries
Светодиод подключается к выводу только через резистор!}}
\end{center}

\subsubsection{Рассчитываем резистор}

Из 3,3~В на светодиод придётся $U_F \approx 2{,}0$~В (для красного), значит, на резистор ---
оставшиеся $3{,}3 - 2{,}0 = 1{,}3$~В. Светодиоду достаточно тока около 6~мА = 0,006~А. По закону Ома:
\[
  R = \frac{U_\text{питания} - U_F}{I} = \frac{3{,}3~\text{В} - 2{,}0~\text{В}}{0{,}006~\text{А}} \approx 217~\si{\ohm}.
\]
Резисторы продаются только определённых номиналов, поэтому берём ближайший: \textbf{\SI{220}{\ohm}}.

Разные светодиоды требуют разного напряжения: красному хватает около 2~В, а синему нужно почти 3~В.%
\footnote{Всё дело в энергии света. Свет состоит из крошечных порций --- фотонов, и у синего света
каждый фотон несёт больше энергии, чем у красного. Чтобы «выбить» из светодиода синий фотон, электрону
нужно больше энергии --- значит, нужно и большее напряжение. Поэтому ультрафиолетовые светодиоды
требуют ещё больше, чем синие, а инфракрасные (в пульте от телевизора) --- меньше всего, около 1,2~В.}
Поэтому и резистор для них нужен разный:

\begin{figure}[H]
\centering
\begin{tikzpicture}
\begin{axis}[
  ybar, width=0.8\textwidth, height=5cm,
  ylabel={прямое напряжение $U_F$, В},
  symbolic x coords={ИК,красный,жёлтый,зелёный,синий,белый},
  xtick=data, x tick label style={font=\sffamily\small},
  tick label style={font=\scriptsize}, label style={font=\small},
  ymin=0, ymax=3.6, bar width=9mm, clip=false,
  nodes near coords, nodes near coords style={font=\scriptsize},
  every node near coord/.append style={/pgf/number format/use comma},
  ymajorgrids, grid style={gray!25},
]
  \addplot[fill=esp!60,draw=espdark] coordinates {(ИК,1.2) (красный,1.9) (жёлтый,2.1) (зелёный,2.2) (синий,3.0) (белый,3.0)};
  \draw[dashed,thick,accent] (rel axis cs:0,0.9167) -- (rel axis cs:1,0.9167) node[pos=0.02,above right,font=\scriptsize,text=accent]{3,3 В --- выход ESP32-S3};
\end{axis}
\end{tikzpicture}
\caption{Прямое напряжение светодиодов разных цветов (примерные значения)}
\end{figure}

\begin{table}[H]
\centering\small
\begin{tabular}{@{}lccc@{}}
\toprule
Цвет светодиода & $U_F$, В & На резисторе, В & Резистор для 6~мА\\
\midrule
Красный & 1,9 & $3{,}3-1{,}9=1{,}4$ & $1{,}4 / 0{,}006 \approx$ \SI{230}{\ohm} → \SI{220}{\ohm}\\
Жёлтый, зелёный & 2,1--2,2 & 1,1--1,2 & $\approx$ \SI{190}{\ohm} → \SI{180}{\ohm}--\SI{220}{\ohm}\\
Синий, белый & 3,0 & 0,3 & $0{,}3 / 0{,}006 = $ \SI{50}{\ohm} → \SI{47}{\ohm}\\
\bottomrule
\end{tabular}
\caption{Пример расчёта резисторов при питании от вывода 3,3~В}
\end{table}

Синему и белому светодиоду от 3,3~В достаётся совсем небольшой запас, поэтому они светят чуть слабее.
Для опытов в курсе лучше всего подходят красные, жёлтые и зелёные светодиоды с резистором \SI{220}{\ohm}.

\subsubsection{Как узнать сопротивление резистора}

Номинал резистора закодирован цветными полосками. Первые две --- цифры, третья --- сколько нулей
дописать, а золотая полоска с краю --- точность:

\begin{figure}[H]
\centering
\begin{tikzpicture}[x=1cm,y=1cm]
  \draw[line width=1.5pt,black!45] (-3.4,0) -- (3.4,0);
  \fill[resbody,draw=black!45,rounded corners=3mm] (-2,-0.55) rectangle (2,0.55);
  \fill[rr] (-1.35,-0.55) rectangle (-1.0,0.55);
  \fill[rr] (-0.7,-0.55) rectangle (-0.35,0.55);
  \fill[rn] (-0.05,-0.55) rectangle (0.3,0.55);
  \fill[rgold] (1.1,-0.55) rectangle (1.45,0.55);
  \foreach \x/\lx/\t in {-1.175/-3.2/{красный: \textbf{2}},-0.525/-1.2/{красный: \textbf{2}},0.125/1.0/{коричневый: \textbf{×10}},1.275/3.2/{золотой: ±5\%}}
    \draw[thin] (\x,-0.6) -- (\x,-0.9) -- (\lx,-1.2) node[below,font=\sffamily\scriptsize,align=center]{\t};
  \node[font=\sffamily\bfseries] at (0,1.1) {2 2 0 = \SI{220}{\ohm}};
\end{tikzpicture}

\vspace{2mm}
\small
\begin{tabular}{@{}*{10}{c}@{}}
\toprule
\colorbox{black}{\textcolor{white}{0}} & \colorbox{rn}{\textcolor{white}{1}} & \colorbox{rr}{\textcolor{white}{2}} & \colorbox{ro}{3} & \colorbox{ry}{4} &
\colorbox{rg}{\textcolor{white}{5}} & \colorbox{rb}{\textcolor{white}{6}} & \colorbox{rv}{\textcolor{white}{7}} & \colorbox{gray}{\textcolor{white}{8}} & \fcolorbox{gray}{white}{9}\\
чёрный & коричн. & красный & оранж. & жёлтый & зелёный & синий & фиолет. & серый & белый\\
\bottomrule
\end{tabular}
\caption{Цветовая маркировка резисторов. Например, коричневый--чёрный--оранжевый = 1, 0 и три нуля = \SI{10000}{\ohm} = \SI{10}{k\ohm}}
\end{figure}

\subsection{Макетная плата и два вида схем}

Чтобы собирать схемы без паяльника, используют \textbf{макетную плату} (англ. \emph{breadboard}). Внутри
под рядами отверстий спрятаны металлические полоски-зажимы. Всё, что воткнуто в одну полоску,
соединено между собой:

\begin{figure}[H]
\centering
\begin{tikzpicture}[bb,x=3mm,y=3mm]
  \bbBoard{24}
  \bbRail{bbred}{0}{24} \bbRail{bbblue}{1}{24}
  \bbRail{bbblue}{16}{24} \bbRail{bbred}{17}{24}
  \foreach \c in {3,4,5}{\bbStrip{okgreen}{\c}{3}{7}}
  \foreach \c in {3,4,5}{\bbStrip{warn}{\c}{10}{14}}
  \bbStrip{okgreen}{15}{3}{7}
  \bbStrip{warn}{15}{10}{14}
  \node[font=\sffamily\scriptsize,anchor=west,align=left] (t1) at (27,17) {шины питания:\\ соединены по всей длине};
  \draw[thin] (t1.west) -- (24.3,16.5);
  \node[font=\sffamily\scriptsize,anchor=west,align=left] (t2) at (27,12) {столбик из 5 отверстий\\ (ряды a--e) соединён};
  \draw[thin] (t2.west) -- (15.3,12);
  \node[font=\sffamily\scriptsize,anchor=west,align=left] (t3) at (27,8.5) {канавка: половинки\\ между собой \textbf{не} соединены};
  \draw[thin] (t3.west) -- (24.8,8.5);
  \node[font=\sffamily\scriptsize,anchor=west,align=left] (t4) at (27,5) {столбик из 5 отверстий\\ (ряды f--j) соединён};
  \draw[thin] (t4.west) -- (15.3,5);
\end{tikzpicture}
\caption{Как соединены отверстия макетной платы (подсвечены соединённые группы)}
\end{figure}

\begin{caution}
Обе ножки детали нельзя втыкать в один столбик: тогда они замкнутся друг на друга, и ток пойдёт мимо детали.
\end{caution}

Каждую схему в этой книге мы рисуем \textbf{двумя способами}:
\begin{itemize}
  \item \textbf{монтажная схема} --- как детали и провода выглядят на макетной плате. По ней удобно собирать;
  \item \textbf{принципиальная электрическая схема} --- детали заменены условными значками, а
        провода --- линиями. По ней удобно понимать, \emph{как схема работает}. Такие схемы рисуют
        инженеры всего мира, и их стоит научиться читать.
\end{itemize}

\begin{table}[H]
\centering\small
\renewcommand{\arraystretch}{2.2}
\begin{tabular}{@{}m{30mm}>{\centering\arraybackslash}m{36mm}>{\centering\arraybackslash}m{40mm}@{}}
\toprule
Деталь & На макетной плате & На принципиальной схеме\\
\midrule
Светодиод & \tikz[bb,x=4mm,y=4mm,baseline=4mm]{\bbLED{red}{0}{0}{2.6}} & \begin{circuitikz}[baseline=-1mm]\draw (0,0) to[led] (2,0);\end{circuitikz}\\
Резистор & \tikz[bb,x=4mm,y=4mm,baseline=-1mm]{\bbR{0}{0}{4}{0}{rr}{rr}{rn}} & \begin{circuitikz}[baseline=-1mm]\draw (0,0) to[R] (2,0);\end{circuitikz}\\
Кнопка & \tikz[bb,x=4mm,y=4mm,baseline=33mm]{\bbButton{0}} & \begin{circuitikz}[baseline=-1mm]\draw (0,0) to[push button] (2,0);\end{circuitikz}\\
Потенциометр & \tikz[bb,x=4mm,y=4mm,baseline=8mm]{\bbPot{0}{0}{10k}} & \begin{circuitikz}[baseline=-1mm]\draw (0,0) to[pR] (2,0);\end{circuitikz}\\
Земля (GND) & черный провод к шине «−» & \begin{circuitikz}[baseline=-3mm]\draw (0,0) node[ground]{};\end{circuitikz}\\
Питание 3,3 В & красный провод к шине «+» & \begin{circuitikz}[baseline=-1mm]\draw (0,-0.3) -- (0,0) node[vcc]{3,3 В};\end{circuitikz}\\
Вывод платы & цветной провод к гребёнке & \begin{circuitikz}[baseline=-1mm]\draw (0,0) node[left]{\pin{4}} to[short,o-] (1,0);\end{circuitikz}\\
\bottomrule
\end{tabular}
\caption{Как одни и те же детали выглядят на двух видах схем}
\end{table}

\begin{tip}
Договоримся о цветах проводов: \textcolor{wred}{\textbf{красный}} --- питание 3,3~В,
\textbf{чёрный} --- земля GND, остальные цвета --- сигналы. В настоящих устройствах инженеры
тоже соблюдают такие правила: так проще искать ошибки.
\end{tip}

\subsection{Подключаем светодиод}

Соберём схему: вывод \pin{4} → резистор \SI{220}{\ohm} → анод светодиода, катод → земля (GND).

\begin{twoview}{Подключение светодиода к выводу GPIO4}{fig:led}
\begin{tikzpicture}[bb]
  \bbBoard{24}
  \bbDevKit{29}{25}
  \bbR{10}{13}{14}{13}{rr}{rr}{rn}
  \bbLED{red}{14}{11}{-1.8}
  \draw[wire=wyellow] (dk-4) -- (10,22) -- (10,14);
  \draw[wire=wblack] (15,14) -- (15,16);
  \draw[wire=wblack] (dk-GND) -- (26.5,4) -- (26.5,16) -- (24,16);
\end{tikzpicture}
\nextview
\begin{circuitikz}
  \draw (0,0) node[left]{\pin{4}} to[short,o-] (0.5,0)
        to[R, l=$R$, a=\SI{220}{\ohm}] (3.3,0)
        to[led, l=LED] (5.8,0) -- (5.8,-1.2) node[ground]{};
  \draw[->,thick,esp] (0.8,0.9) -- node[above,font=\sffamily\small]{ток $I$} (2.8,0.9);
\end{circuitikz}
\end{twoview}

Сравни две картинки: жёлтый провод на макетной плате --- это линия от \pin{4} на схеме, чёрные
провода --- значок земли. Проследи путь тока пальцем: от вывода платы через резистор и светодиод
к земле и обратно в плату --- цепь замкнута.

\subsection{Программа Blink: мигаем светодиодом}

\codecap{Мигание светодиодом}
\codeinput{l05_blink}

\begin{itemize}
  \item \code{pinMode(LED\_PIN, OUTPUT)} --- сделать вывод выходом;
  \item \code{digitalWrite(LED\_PIN, HIGH)} --- выдать логическую 1 (3,3~В), светодиод загорается;
  \item \code{digitalWrite(LED\_PIN, LOW)} --- выдать 0 (0~В), светодиод гаснет;
  \item \code{delay(500)} --- подождать полсекунды.
\end{itemize}

Посмотрим, какое напряжение при этом на выводе. Это цифровой сигнал из урока~\ref{les:signals}:
у него всего два уровня, а переходы между ними резкие. Такой сигнал называют \textbf{прямоугольным}
или меандром:

\begin{figure}[H]
\centering
\begin{tikzpicture}[x=1.15cm,y=1cm]
  \draw[-{Stealth[length=2mm]}] (0,0) -- (8.6,0) node[right]{$t$, мс};
  \draw[-{Stealth[length=2mm]}] (0,0) -- (0,2.2) node[above]{$U$, В};
  \draw[very thick,esp] (0,0) -- (0,1.2) -- (2,1.2) -- (2,0) -- (4,0) -- (4,1.2) -- (6,1.2) -- (6,0) -- (8,0) -- (8,1.2);
  \foreach \x/\l in {0/0,2/500,4/1000,6/1500,8/2000} \draw (\x,0.05)--(\x,-0.05) node[below,font=\scriptsize]{\l};
  \draw (0.05,1.2) -- (-0.05,1.2) node[left,font=\scriptsize]{3,3};
  \node[font=\sffamily\scriptsize,text=esp] at (1,0.6) {горит};
  \node[font=\sffamily\scriptsize,text=gray] at (3,0.3) {не горит};
  \node[font=\sffamily\scriptsize,text=esp] at (5,0.6) {горит};
  \node[font=\sffamily\scriptsize,text=gray] at (7,0.3) {не горит};
  \draw[decorate,decoration={brace,amplitude=4pt}] (0,1.35) -- (2,1.35) node[midway,above=4pt,font=\scriptsize]{\code{HIGH}, 500 мс};
  \draw[decorate,decoration={brace,amplitude=4pt}] (2,1.35) -- (4,1.35) node[midway,above=4pt,font=\scriptsize]{\code{LOW}, 500 мс};
  \draw[{Stealth[length=2mm]}-{Stealth[length=2mm]},accent,thick] (0,-0.75) -- node[below,font=\scriptsize]{период $T$ = 1000 мс = 1 с} (4,-0.75);
\end{tikzpicture}
\caption{Напряжение на выводе \pin{4} при работе Blink --- прямоугольный сигнал}
\end{figure}

Время одного полного цикла «горит + не горит» называется \textbf{периодом} $T$. Здесь
$T = 500 + 500 = 1000$~мс = 1~с. Сколько периодов умещается в секунду, называется \textbf{частотой}:
\[
  f = \frac{1}{T} = \frac{1}{1~\text{с}} = 1~\text{Гц}.
\]
Светодиод мигает один раз в секунду --- частота 1~Гц.

\subsubsection{Что видит глаз}

При частоте 1~Гц глаз прекрасно видит каждую вспышку: светодиод полсекунды горит, полсекунды
не горит. Глаз успевает заметить изменения, потому что они медленные. А что будет, если уменьшать
\code{delay()} и мигать всё быстрее? Светодиод будет мигать чаще и чаще\ldots\ а потом случится
кое-что удивительное. Этот опыт мы проведём в уроке~\ref{les:pwm}, и он приведёт нас к способу
менять яркость светодиода.

\subsection{Мигание без \texttt{delay()}}

У \code{delay()} есть недостаток: пока идёт ожидание, программа \emph{ничего больше не делает}.
Она не заметит нажатия кнопки и не ответит компьютеру. Представь, что ты варишь яйцо и полторы
минуты смотришь на часы, не отходя от плиты. Удобнее иногда поглядывать на часы, а в промежутках
заниматься другими делами. Так же поступает программа ниже: она постоянно сравнивает текущее время
\code{millis()} с моментом последнего переключения.

\begin{figure}[H]
\centering
\begin{tikzpicture}[x=1cm,y=1cm]
  \node[font=\sffamily\small,anchor=east] at (-0.2,1) {с \code{delay()}};
  \fill[esp!35] (0,0.8) rectangle (0.3,1.2);
  \fill[gray!25] (0.3,0.8) rectangle (5,1.2) node[midway,font=\sffamily\scriptsize]{ждём и ничего не делаем};
  \fill[esp!35] (5,0.8) rectangle (5.3,1.2);
  \fill[gray!25] (5.3,0.8) rectangle (10,1.2) node[midway,font=\sffamily\scriptsize]{ждём и ничего не делаем};
  \node[font=\sffamily\small,anchor=east] at (-0.2,0) {с \code{millis()}};
  \foreach \x in {0,0.4,...,9.7} \fill[accent!40] (\x,-0.2) rectangle (\x+0.3,0.2);
  \fill[esp!60] (0,-0.2) rectangle (0.3,0.2);
  \fill[esp!60] (4.8,-0.2) rectangle (5.1,0.2);
  \fill[esp!60] (9.6,-0.2) rectangle (9.9,0.2);
  \node[font=\sffamily\scriptsize,align=center] at (5,-0.8) {\textcolor{accent}{■} много коротких проходов \code{loop()}: «пора?» --- «нет, занимаюсь другим»\quad
     \textcolor{esp!60}{■} «пора!» --- переключаем светодиод};
\end{tikzpicture}
\caption{Слева направо идёт время. Программа с \code{millis()} постоянно свободна}
\end{figure}

\codecap{Мигание без остановки программы}
\codeinput{l05_blink_millis}

Запомни этот приём: \code{if (millis() - last >= INTERVAL)}. Он будет встречаться почти в каждом
уроке.

\subsection{Практика: бегущий огонь}

Подключим четыре светодиода к выводам \pin{4}--\pin{7}, каждый со своим резистором, и будем
зажигать их по очереди. Получится огонёк, который бежит по кругу, как на новогодней гирлянде.

\begin{twoview}{Четыре светодиода на выводах GPIO4--GPIO7}{fig:led4}
\begin{tikzpicture}[bb]
  \bbBoard{24}
  \bbDevKit{29}{25}
  \foreach \k/\pin/\col/\wc in {0/4/red/wyellow,1/5/yellow/wgreen,2/6/green/wblue,3/7/blue/wpurple}{
    \pgfmathtruncatemacro{\c}{1+6*\k}
    \bbR{\c}{13}{\c+4}{13}{rr}{rr}{rn}
    \bbLED{\col}{\c+4}{11}{-1.8}
    \draw[wire=wblack] (\c+5,14) -- (\c+5,16);
    \draw[wire=\wc] (dk-\pin) -- (\c,22-\k) -- (\c,14);
  }
  \draw[wire=wblack] (dk-GND) -- (26.5,4) -- (26.5,16) -- (24,16);
\end{tikzpicture}
\nextview
\begin{circuitikz}[scale=0.9,transform shape]
  \foreach \k/\p in {0/4,1/5,2/6,3/7}{
    \draw (0,-\k*1.3) node[left]{\pin{\p}} to[short,o-] (0.5,-\k*1.3)
          to[R, l=\SI{220}{\ohm}] (3,-\k*1.3) to[led] (5.2,-\k*1.3) -- (6,-\k*1.3);
  }
  \draw (6,0) -- (6,-4.6) node[ground]{};
\end{circuitikz}
\end{twoview}

Обрати внимание: на принципиальной схеме земли всех светодиодов соединены одной линией. На
макетной плате это сделано с помощью шины «−»: все катоды подключены к ней, а она --- к GND платы.

\codecap{Бегущий огонь}
\codeinput{l05_running}

Номера выводов собраны в \textbf{массив} \code{LEDS[]} --- список значений под одним именем.
\code{LEDS[0]} --- это 4, \code{LEDS[1]} --- 5 и так далее. Выражение \code{(current + 1) \% COUNT}
даёт по очереди 1, 2, 3, 0, 1, 2\ldots: знак \code{\%} означает остаток от деления, и после
последнего светодиода счёт начинается снова с нулевого.

\subsection{Встроенный RGB-светодиод}

На плате уже есть светодиод, который светится любым цветом. Внутри него три крошечных светодиода:
красный (R), зелёный (G) и синий (B). Смешивая их свет в разных пропорциях, можно получить любой цвет.
Точно так же устроен каждый пиксель экрана телефона и телевизора!

\begin{figure}[H]
\centering
\begin{tikzpicture}
  \fill[black,rounded corners=2mm] (-2.4,-2.2) rectangle (2.4,2.2);
  \begin{scope}[blend group=screen]
    \fill[red]   (90:0.75) circle (1.15);
    \fill[green] (210:0.75) circle (1.15);
    \fill[blue]  (330:0.75) circle (1.15);
  \end{scope}
  \node[font=\sffamily\scriptsize\bfseries,white] at (90:1.35) {R};
  \node[font=\sffamily\scriptsize\bfseries,white] at (210:1.35) {G};
  \node[font=\sffamily\scriptsize\bfseries,white] at (330:1.35) {B};
  \node[font=\sffamily\small,anchor=west,align=left] at (2.8,0) {красный + зелёный = \textbf{жёлтый}\\
     зелёный + синий = \textbf{голубой}\\ красный + синий = \textbf{пурпурный}\\ все три = \textbf{белый}};
\end{tikzpicture}
\caption{Смешение цветов света. Это не то же самое, что смешивать краски: краски вместе дают тёмный цвет, а свет --- белый}
\end{figure}

Встроенный светодиод подключён к \pin{48} (на некоторых платах --- к \pin{38}). Для него есть готовая
функция \code{rgbLedWrite(pin, r, g, b)}, где каждая яркость --- от 0 до 255:

\codecap{Смена цветов встроенного RGB-светодиода}
\codeinput{l05_rgb}

\begin{summary}
\begin{itemize}[leftmargin=*]
  \item Светодиод пропускает ток в одну сторону: длинная ножка --- анод (+), короткая --- катод (−).
  \item Светодиод подключается только через резистор. Резистор считаем по закону Ома: $R = (U - U_F)/I$.
  \item Чем «синее» светодиод, тем больше напряжения он требует.
  \item Blink создаёт на выводе прямоугольный сигнал. Период $T$ --- время одного цикла, частота $f = 1/T$.
  \item Чтобы программа не «засыпала», время отсчитываем через \code{millis()}, а не \code{delay()}.
\end{itemize}
\end{summary}

\begin{quiz}
\begin{enumerate}
  \item Что случится, если подключить светодиод к выводу без резистора?
  \item Какой резистор нужен зелёному светодиоду ($U_F = 2{,}2$~В) при токе 5~мА и питании 3,3~В?
  \item Светодиод горит 200~мс и не горит 300~мс. Чему равны период и частота?
  \item Какие цвета полосок у резистора \SI{4.7}{k\ohm} (4700~\si{\ohm})?
\end{enumerate}
\end{quiz}

\begin{tasks}
\begin{enumerate}
  \item Собери «светофор» из красного, жёлтого и зелёного светодиодов с правильной последовательностью и длительностями.
  \item Сделай два светодиода, мигающих с разной частотой (1~Гц и 3~Гц), без \code{delay()}.
  \item Заставь «бегущий огонь» бегать туда и обратно, как фары автомобиля из старого фильма.
  \item Сделай RGB-светодиод «радугой»: красный → жёлтый → зелёный → голубой → синий → пурпурный.
\end{enumerate}
\end{tasks}
